\scanpage{4}
L'assertion précédant la proposition est équivalente aux suivantes
\[
\begin{aligned}
D+N&=\mathcal E_0\\
D\cap N&=L\\[4pt]
\Phi+N&=\mathcal E\qquad\deleted{\unsure{annulation}}\\
\Phi\cap N&=\Phi_N
\end{aligned}
\]

Définition de l'application $\mathcal E\longrightarrow \unsure{Z|\mathbb Z}+\unsure{\mathfrak P|\mathfrak D}$

sur $D$ : coïncide avec l'application \uline{canonique} $D\mathrel{\text{---}}\unsure{\mathfrak P|\mathfrak D}$

sur $\Phi$ : Soit $E$ un fibré vectoriel, fibre type $K^n$ et de classe de Chern
$c^1(E)\in\unsure{\mathfrak P|\mathfrak D}$, on prend
\[
h(E)=(n,c^1(E))
\]

Soit $F$ un faisceau algébrique cohérent, $c(F)\in \unsure{Z|\mathbb Z}+\unsure{\mathfrak P|\mathfrak D}$,
\unsure{je} dis que
\[
(1)\qquad \unsure{\chi|\xi}(F)=K_1\bigl(c(F).(1-\tfrac12 K)\bigr)
\qquad (K\in\unsure{\mathfrak P|\mathfrak D})
\]
où $K$ est la classe canonique ; $K_1$ est la forme linéaire sur
$\unsure{Z|\mathbb Z}+\unsure{\mathfrak P|\mathfrak D}$ égale à $0$ sur $\unsure{Z|\mathbb Z}$, et sur
$\unsure{\mathfrak P|\mathfrak D}$ : l'application ``degré''
$\unsure{\mathfrak P|\mathfrak D}\to \unsure{Z|\mathbb Z}$.
Pour vérifier (1), il suffit de la vérifier pour $F=k_x$ et pour $F=\underline O_x$

\diagram{\begin{tikzcd}[row sep=large,column sep=large,
execute at end picture={
\draw[decorate,decoration={brace,amplitude=5pt}]
(d.east) -- (e.north west) node[midway,above right] {$\Phi$};
\draw[decorate,decoration={brace,mirror,amplitude=5pt}]
(n.south east) -- (e.south) node[midway,below right,align=center] {$\mathcal E/N$ classes de\\faisceaux};
}]
& |[alias=d]| D \arrow[dr] & & \\
L=D\cap N \arrow[ur,"\unsure{\mathfrak P|\mathfrak D}"] \arrow[dr,"?"'] & & D+N=\mathcal E_0 \arrow[r,"\mathbb Z"'] & |[alias=e]| \mathcal E\\
& |[alias=n]| N \arrow[ur] & &
\end{tikzcd}}

\margin{\unsure{LINEN TYPEWRITER}\quad
\figure{\begin{tikzpicture}[scale=.13]
\foreach \a in {0,45,...,315}{\draw (\a:1) circle (.54);}
\draw (0,0) circle (.36);
\end{tikzpicture}}\quad \unsure{ARGENTINA}}

\margin{3}
